\section{Introduction}
\label{sec:intro}
\guangyao{I feel that you can directly use the narrative from MemRL: Self-Evolving Agents via Runtime Reinforcement Learning on Episodic Memory if the following questions can be answered: 1. why should we consider LLM agent for video game? is this well-justified in the existing literature or somewhere? 2. if 1 is justified, evolving memory makes sense easily because we want it to improve without changing the backbone LLM. But the next question is how is video game fundamentally different langugage-Q/A type tasks? 3. how do differences mentioned in 2 motivate new staff in your approach?  One point occurs to me is that the video game may have some period pattern, then somehow you exploit such a feature to design a high-efficient somehow skill selection and evolving mechanism.}

% --- (1) Why LLM agents for video games? (literature-justified) ---
Video games are a well-established testbed for evaluating and advancing LLM and VLM agents~\cite{berner2019dota, xu2025agents}. Prior work argues that gameplay demands the same core capabilities as general-purpose agents---long-context understanding, memory, long-horizon planning, and adaptive decision-making---in scalable, controllable, and interactive environments~\cite{wang2023voyager, raad2024scaling}, with direct implications for continual-learning and deployable agents in real-world settings~\cite{hu2024survey}. Games are thus not only a challenging benchmark but a justified target for agent research: they stress capabilities that transfer beyond narrow Q/A. At the same time, learning from interaction in games is difficult~\cite{liao2025think, lu2025cultivating}: rewards are delayed~\cite{liu2025agentic}, strategies are compositional, and high-quality human demonstrations or skill annotations are scarce. Many existing game-playing agents therefore still rely on curated human data~\cite{hu2023language} or prolonged training with human or external-LLM feedback.

% --- (2) Evolving memory without changing backbone; (3) How games differ from Q/A and motivate our approach (periodic patterns → skill selection & evolution) ---
Recent progress in self-improving and self-evolving LLM agents~\cite{li2025self, wang2025vision, he2025visplay, huang2025r} suggests a path beyond annotation-heavy pipelines: agents can improve from large amounts of unlabeled interaction data and reduce reliance on high-quality demonstrations~\cite{gao2025survey, xiangsystematic}. In particular, \emph{evolving memory}---updating experience and reusable structures at runtime without changing the backbone LLM---allows the agent to get better over time while keeping the same model. This is especially attractive for games, where collecting and curating data is costly. A natural next question is: \textbf{how do video games differ from language Q/A or shorter-horizon tasks, and how should the agent's memory and improvement mechanism be designed accordingly?}

Unlike single-turn or short-horizon Q/A, video game episodes are \textbf{long-horizon} and \textbf{multi-hop}: a single episode requires the agent to sequentially deploy and coordinate \emph{multiple} distinct skills (e.g., clear rows, merge tiles, avoid loss), each spanning many turns. Moreover, gameplay exhibits \textbf{recurring, periodic structure}: the same types of subgoals and skill-level behaviors repeat across episodes and across games (e.g., navigation, clearing, survival). Treating an episode as a flat trajectory ignores this structure and hinders experience reuse---useful behaviors are temporally entangled, context-dependent, and hard to retrieve or improve in isolation. These differences motivate two design choices. \textbf{First}, we decompose long trajectories into \textbf{skill-level units} via unsupervised segmentation, so that the agent can store, retrieve, and refine \emph{reusable} behaviors rather than raw sequences. \textbf{Second}, we exploit the recurring nature of subgoals to build an \textbf{agent-managed skill bank} that is both \emph{queryable} (efficient skill selection via retrieval and applicability scoring) and \emph{evolvable} (continual refinement from new rollouts). Memory and skills thus evolve at runtime without changing the backbone LLM~\cite{xu2025mem, sarch2024vlm}, and recent work shows that such structured memory and skill management can improve performance on complex multi-hop tasks in web and embodied agents~\cite{xia2026skillrl, zhang2026memrl, zhang2026memskill}. Despite this, skill segmentation and skill-bank evolution for LLM agents in long-horizon gameplay remain underexplored; changing game states, branching outcomes, and sparse feedback make the problem substantially harder than in shorter or more structured settings~\cite{berner2019dota, raad2024scaling}.

% --- Our framework ---
In this paper, we propose a \textbf{multi-agent co-evolution} framework for multi-hop game playing that closes the loop between (i)~a language \textbf{decision agent} that plays games using primitive actions and tool calls (e.g., \texttt{QUERY\_SKILL}, \texttt{CALL\_SKILL}, \texttt{QUERY\_MEM}) and (ii)~an \textbf{agent-managed skill bank} that performs unsupervised skill discovery, maintenance, and retrieval. The decision agent consumes the skill bank for retrieval-augmented planning (protocols) and reward shaping (effect contracts); the skill bank agent ingests the decision agent's trajectories and converts them into segmented, labeled skills with protocols and contracts via a multi-stage pipeline (boundary proposal, segment decoding, contract learning, bank maintenance). Both agents are trained in the loop: the decision agent with GRPO and the skill bank with Hard-EM and LoRA adapters, so that better rollouts yield better skills and a better skill bank yields better decision-making. To the best of our knowledge, this is the first framework to jointly couple LLM-based game decision-making with an agentic skill-bank pipeline for unsupervised multi-hop skill segmentation and continual skill refinement in a unified co-evolution loop.

Our main contributions include:
\begin{itemize}
    \item We present a general co-evolution framework that closes the loop between gameplay and skill acquisition: an LLM decision agent interacts with the environment to collect trajectories, and an agent-managed skill pipeline converts them into reusable skills (protocols and effect contracts) that improve future gameplay via retrieval and reward shaping.
    \item We propose a practical agentic pipeline for skill bank maintenance that transforms unlabeled, multi-hop trajectories into reusable skills via segmentation, contract learning, and bank updates. We exploit the recurring, periodic structure of gameplay to design an efficient skill-selection mechanism (RAG-based retrieval with applicability scoring) and an evolving skill bank that is refined from rollouts without changing the backbone LLM.
    \item We conduct extensive experiments across multiple video game environments requiring multi-hop skill usage, showing that the proposed co-evolution framework outperforms vanilla LLM agents, non-co-evolving methods, and training-free baselines. \textcolor{red}{[show results]}
\end{itemize}
